\documentclass[11pt]{article}
\usepackage[margin=1in]{geometry}
\usepackage{amsmath,amssymb,booktabs,tabularx,listings,xcolor,hyperref}
\hypersetup{colorlinks=true,linkcolor=blue!50!black}
\lstset{basicstyle=\ttfamily\small,columns=fullflexible,keepspaces=true,breaklines=true,frame=single}
\setlength{\parindent}{0pt}
\setlength{\parskip}{6pt}
\setlength{\emergencystretch}{2em}
\title{CSI 2110 --- Lecture 2\\Stacks, Queues, and Deques}
\author{Study notes based on Bo Sun's Fall 2026 lecture}
\date{}
\begin{document}
\maketitle
\textbf{Source:} \texttt{lec02-stacks, queues, deques.pdf}, 53 PDF pages. References use PDF page numbers. Pseudocode is expanded to make boundary cases and updates explicit; complexity explanations and review answers supplement the slides.

\section{Abstract data types}
\textit{Source: pp.\ 4--5.}

An \textbf{abstract data type (ADT)} specifies what values can be stored and what operations do. A \textbf{concrete data structure} specifies how those operations are implemented. In Java, an interface can express the ADT and a class can implement it.

The lecture's stack, queue, and deque interfaces return \texttt{null} for an empty removal or empty endpoint inspection. Fixed-capacity insertion into a full array reports an error. These are the course conventions; they are not a claim about every Java library method.

\section{Stacks: last in, first out}
\textit{Source: pp.\ 7--19.}

A stack adds and removes at the \textbf{top}. The most recently pushed element is the next one popped.
\begin{center}
\begin{tabularx}{\textwidth}{lX}
\toprule
Operation & Meaning\\
\midrule
\texttt{push(e)} & Put $e$ on top.\\
\texttt{pop()} & Remove and return the top; return null if empty.\\
\texttt{top()} & Inspect the top without removing it; null if empty.\\
\texttt{size()}, \texttt{isEmpty()} & Return the size or test whether it is zero.\\
\bottomrule
\end{tabularx}
\end{center}
Example: typing C, A, T and undoing twice removes T and then A, leaving C.

\subsection{Array stack}
Use an array $S$ of capacity $N>0$ and a top index $t$, initially $-1$. The occupied segment is $S[0\ldots t]$, and size is $t+1$.
\begin{lstlisting}
push(e):
    if t+1 = N: error "full stack"
    t <- t+1
    S[t] <- e

pop():
    if t = -1: return null
    result <- S[t]
    S[t] <- null
    t <- t-1
    return result

top():
    if t = -1: return null
    return S[t]
\end{lstlisting}
All operations are $O(1)$, and storage is $O(N)$. Clearing a removed reference prevents the array from retaining the object unnecessarily. A resizable array changes the worst-case cost of a push that triggers copying; Lecture 3 explains amortized bounds.

\subsection{Singly linked stack}
Make the head of a singly linked list the stack top; maintain a size counter.
\begin{lstlisting}
push(e):
    head <- new Node(e, head)
    size <- size+1

pop():
    if head = null: return null
    result <- head.element
    head <- head.next
    size <- size-1
    return result
\end{lstlisting}
All stack operations are $O(1)$ and storage is $O(n)$. Using the tail of a singly linked list as the top would make removal expensive.

\section{Stack applications}
\textit{Source: pp.\ 20--25.}

\subsection{Reverse an array}
Push the elements left to right, then pop them back into the array left to right. The input $[4,1,2,3]$ becomes $[3,2,1,4]$. Time is $\Theta(n)$ and auxiliary stack space is $\Theta(n)$. In Java, a \texttt{Stack<Integer>} example should use an \texttt{Integer[]} argument when the method signature requires that array type; \texttt{int[]} is a different type.

\subsection{Match delimiters}
Push each opening delimiter. For a closing delimiter, reject if the stack is empty or its popped opening delimiter has a different type. Ignore non-delimiter characters. Accept at the end only if the stack is empty.
\begin{lstlisting}
isMatched(text):
    S <- empty stack
    for each character c in text:
        if c is one of '(', '[', '{':
            S.push(c)
        else if c is one of ')', ']', '}':
            if S.isEmpty(): return false
            if S.pop() does not match c: return false
    return S.isEmpty()
\end{lstlisting}
The stack contains the unmatched opening delimiters in nesting order. A closing delimiter must match the \emph{most recent} unmatched opening. Thus \texttt{([]\{\})} is valid, while \texttt{([)]}, \texttt{(()}, and \texttt{)(} are invalid. Time is $O(n)$ and worst-case space is $O(n)$.

\subsection{Evaluate arithmetic expressions}
Maintain a value stack $V$ and an operator stack $O$. For the left-associative operators in the lecture, before pushing a new operator, apply pending operators of higher or equal precedence. At the end, apply all remaining operators.

To apply a binary operator, pop the \textbf{right operand first}, then the left, and push $\text{left}\ \text{operator}\ \text{right}$. This matters for subtraction and division.
\[
14-3\times2+7\;\longrightarrow\;14-6+7
\;\longrightarrow\;8+7\;\longrightarrow\;15.
\]
When the $+$ arrives, apply $\times$ and then $-$ before storing $+$. A parenthesis-aware extension pushes opening parentheses and applies operators until that opening marker when a closing parenthesis arrives. Equal-precedence popping assumes left associativity; it must be adjusted for any right-associative operators.

\section{Queues: first in, first out}
\textit{Source: pp.\ 27--39.}

A queue inserts at the \textbf{rear} and removes at the \textbf{front}. \texttt{enqueue(e)} inserts; \texttt{dequeue()} removes and returns the oldest element; \texttt{first()} inspects it. The ADT also includes \texttt{size()} and \texttt{isEmpty()}. A shared printer is a typical arrival-order application.

\subsection{Circular array queue}
Maintain capacity $N>0$, front index $f$, and size $s$, initially $f=s=0$. Wrap indices using modulo:
\[
\text{next insertion index }r=(f+s)\bmod N,
\qquad\text{last occupied index }(f+s-1)\bmod N.
\]
The last-index formula is used only when $s>0$. Logical element $i$ is stored at $(f+i)\bmod N$.
\begin{lstlisting}
enqueue(e):
    if s = N: error "full queue"
    Q[(f+s) mod N] <- e
    s <- s+1

dequeue():
    if s = 0: return null
    result <- Q[f]
    Q[f] <- null
    f <- (f+1) mod N
    s <- s-1
    return result
\end{lstlisting}
All operations are $O(1)$; storage is $O(N)$. No element shifting is needed. \textbf{$f=r$ alone does not distinguish empty from full}: it occurs when $s=0$ and when $s=N$. The size counter resolves the ambiguity.

\textbf{Wraparound trace from the slides.} With $N=11$, $f=7$, and logical contents $[A,B,C]$, the next insertion slot is $10$. Enqueue D at $10$, then E at $0$. Dequeue returns A, sets $f=8$, and leaves $s=4$ with logical contents $[B,C,D,E]$.

\subsection{Singly linked queue}
Use the head as the front and tail as the rear. Enqueue by appending at the tail; dequeue by removing the head. Maintain the size. When adding to an empty queue, set both pointers. When the last element is removed, set the tail to null as well.

This gives $O(1)$ operations and $O(n)$ storage. Reversing the endpoint roles would require removing the tail of a singly linked list, taking $O(n)$ time. An array queue is useful when capacity is bounded; a linked or extendable-array implementation supports growth.

\section{Queue application: recent calls}
\textit{Source: pp.\ 40--41.}
Calls arrive at strictly increasing timestamps $t$. Return how many lie in the inclusive interval $[t-3000,t]$:
\begin{lstlisting}
ping(t):
    Q.enqueue(t)
    cutoff <- t-3000
    while not Q.isEmpty() and Q.first() < cutoff:
        Q.dequeue()
    return Q.size()
\end{lstlisting}
The comparison is strict: a timestamp equal to the cutoff remains. Calls at $10,500,3010,3501$ produce counts $1,2,3,2$. At $3501$, the cutoff is $501$, so $10$ and $500$ are removed.

\textbf{Additional analysis:} one call can remove many timestamps, but each timestamp is enqueued once and dequeued at most once. Across $m$ calls the queue work is $O(m)$, or $O(1)$ amortized per call. Storage is proportional to the number of calls still in the time window.

\section{Deques: both endpoints}
\textit{Source: pp.\ 43--51.}

A \textbf{double-ended queue} supports \texttt{addFirst}, \texttt{addLast}, \texttt{removeFirst}, and \texttt{removeLast}. Inspectors are \texttt{first}, \texttt{last}, \texttt{size}, and \texttt{isEmpty}.

\begin{center}
\begin{tabular}{lll}
\toprule
Use a deque as & Insert & Remove\\
\midrule
Stack, top at back & \texttt{addLast(e)} & \texttt{removeLast()}\\
Queue & \texttt{addLast(e)} & \texttt{removeFirst()}\\
\bottomrule
\end{tabular}
\end{center}

\subsection{Circular array implementation}
Use the same $f,s,N$ invariant as a queue. \texttt{addLast} and \texttt{removeFirst} are the queue operations. The two remaining updates are:
\begin{lstlisting}
addFirst(e):
    if s = N: error "full deque"
    f <- (f-1+N) mod N
    D[f] <- e
    s <- s+1

removeLast():
    if s = 0: return null
    last <- (f+s-1) mod N
    result <- D[last]
    D[last] <- null
    s <- s-1
    return result
\end{lstlisting}
Adding $N$ before taking the remainder in \texttt{addFirst} avoids a negative index in languages such as Java. All operations are $O(1)$ with a fixed-capacity array; space is $O(N)$.

\subsection{Doubly linked implementation}
Use header and trailer sentinel nodes. The first element is \texttt{header.next}; the last is \texttt{trailer.prev}. Inserting or removing next to a sentinel requires only local link changes. All operations are $O(1)$ and space is $O(n)$. A singly linked list cannot supply all four endpoint updates in constant time.

\subsection{Palindrome application}
Put the characters into a deque. While at least two remain, remove and compare the first and last. Reject a mismatch; otherwise accept when zero or one remains. For example, RADAR matches R/R, then A/A, leaving D. Time is $O(n)$ and deque space is $O(n)$. Case folding and punctuation removal, if desired, are separate preprocessing choices.

\section{Worked exercise and checks}
\textit{Source exercise: p.\ 52.}
Push and enqueue $0,1,2,3,4$ into an empty stack and queue. Then pop four stack elements and enqueue them. Finally, dequeue twice and push those values onto the stack.
\begin{center}
\begin{tabular}{lll}
\toprule
Checkpoint & Stack (bottom to top) & Queue (front to rear)\\
\midrule
After initial loop & $[0,1,2,3,4]$ & $[0,1,2,3,4]$\\
After four transfers & $[0]$ & $[0,1,2,3,4,4,3,2,1]$\\
After two final pushes & $[0,0,1]$ & $[2,3,4,4,3,2,1]$\\
\bottomrule
\end{tabular}
\end{center}
Check that you can explain why a stack reverses order, why a queue preserves it, why circular arrays need both an index and a size, and why a linked deque needs backward links for efficient tail deletion.
\end{document}
